% Original complete parameter shift and finite shot worked derivation.
\subsection*{Worked derivation: parameter shift and the measurement cost}
The parameter-shift formula above is simple to quote.
To see why it is exact, start from a single rotation
$\op R(\theta)=\exp(-\ii\theta\sigma_\mu/2)$.
Since $\sigma_\mu^2=\mathbf1$, even and odd powers in the
exponential sum separately to
\begin{equation}
\op R(\theta)=\mathbf1\cos\frac\theta2
-\ii\sigma_\mu\sin\frac\theta2.
\end{equation}
For any fixed input density matrix $\rho$ and Hermitian
measurement operator $\op O$, write
$f(\theta)=\operatorname{tr}[\op O\op R(\theta)
\rho\op R^\dagger(\theta)]$. Substitute the two-term
rotation decomposition and expand the four products.
Terms containing two factors $\cos(\theta/2)$ or two
factors $\sin(\theta/2)$ give
$\cos^2(\theta/2)$ and $\sin^2(\theta/2)$;
the two cross terms give
$\sin(\theta/2)\cos(\theta/2)$.
Using
\begin{equation}
\cos^2\frac\theta2=\frac{1+\cos\theta}2,\qquad
\sin^2\frac\theta2=\frac{1-\cos\theta}2,\qquad
2\sin\frac\theta2\cos\frac\theta2=\sin\theta,
\end{equation}
we obtain
\begin{equation}
f(\theta)=A+B\cos\theta+C\sin\theta,
\end{equation}
for real coefficients $A,B,C$ independent of $\theta$.
This holds with arbitrary fixed unitaries before and after
the displayed rotation; they simply modify the effective
$\rho$ and $\op O$.
Differentiate analytically:
\begin{equation}
\frac{\dd f}{\dd\theta}=-B\sin\theta+C\cos\theta.
\end{equation}
Now evaluate the function at shifted arguments:
\begin{align}
f(\theta+\pi/2)&=A-B\sin\theta+C\cos\theta,\\
f(\theta-\pi/2)&=A+B\sin\theta-C\cos\theta.
\end{align}
Subtracting and dividing by two gives
\begin{equation}
\boxed{f'(\theta)=\frac12
\left[f(\theta+\pi/2)-f(\theta-\pi/2)\right],}
\end{equation}
which is exactly Eq.~\eqref{eq:qnn-shift}.

\paragraph{The gradient is exact; measurements introduce noise.}
For each of the two shifted circuits, suppose we measure
a Pauli observable with independent outputs $z=\pm1$.
If the true expectation is $f_\pm$, then
$E[z]=f_\pm$, $E[z^2]=1$ and
$\operatorname{Var}(z)=1-f_\pm^2$.
With $M$ independent shots per shift,
the sample mean $\widehat f_\pm$ has variance
\begin{equation}
\operatorname{Var}(\widehat f_\pm)
=\frac{1-f_\pm^2}{M}.
\end{equation}
The two batches are independent, so for
$\widehat g=(\widehat f_+-\widehat f_-)/2$,
\begin{align}
E[\widehat g]
&=\tfrac12(f_+-f_-)=f'(\theta),\\
\operatorname{Var}(\widehat g)
&=\tfrac14\bigl[\operatorname{Var}(\widehat f_+)
+\operatorname{Var}(\widehat f_-)\bigr]\notag\\
&=\boxed{\frac{2-f_+^2-f_-^2}{4M}
\leq\frac1{2M}.}
\end{align}
The estimator is unbiased, and its uncertainty depends on
measurement outcomes as well as the shot budget.

\paragraph{A useful counterexample to a universal claim.}
Consider $n$ independently rotated qubits
$\ket{\psi(\vect\theta)}=\bigotimes_{j=1}^n
R_y(\theta_j)\ket0$. Measuring the global observable
$Z^{\otimes n}$ gives
\begin{equation}
f_{\mathrm{global}}(\vect\theta)
=\prod_{j=1}^n\cos\theta_j,\qquad
\partial_{\theta_1}f_{\mathrm{global}}
=-\sin\theta_1\prod_{j=2}^n\cos\theta_j.
\end{equation}
When each $\theta_j$ is independent and uniform
on $[-\pi,\pi]$, the averages
$E[\sin^2\theta_j]=E[\cos^2\theta_j]=1/2$
give
\begin{equation}
E[(\partial_{\theta_1}f_{\mathrm{global}})^2]
=\left(\frac12\right)^n.
\end{equation}
For the local observable $Z_1$, however,
$f_{\mathrm{local}}=\cos\theta_1$ and
\begin{equation}
E[(\partial_{\theta_1}f_{\mathrm{local}})^2]
=\frac12
\end{equation}
for \emph{every} $n$. The same product circuit can
therefore have exponentially small global gradients and
a dimension-independent local gradient. The precise
cost observable and parameter ensemble matter.
The
\href{https://kangqiaoliu.github.io/courses/advanced-quantum-mechanics/companion/labs/IC-09-06/index.html}
{gradient concentration lab} tests these exact moments
against reproducible random ensembles.

\paragraph{Checkpoint.}
At $n=10$, calculate the expected squared global
gradient and compare it with the local value.
For $M=512$, compute the conservative standard-deviation
bound $\sqrt{1/(2M)}$ for a two-shift estimator.
Which gradients are easier to distinguish from shot noise?
